% Options for packages loaded elsewhere
\PassOptionsToPackage{unicode}{hyperref}
\PassOptionsToPackage{hyphens}{url}
\PassOptionsToPackage{dvipsnames,svgnames,x11names}{xcolor}
%
\documentclass[
  11pt,
]{article}
\usepackage{amsmath,amssymb}
\usepackage{iftex}
\ifPDFTeX
  \usepackage[T1]{fontenc}
  \usepackage[utf8]{inputenc}
  \usepackage{textcomp} % provide euro and other symbols
\else % if luatex or xetex
  \usepackage{unicode-math} % this also loads fontspec
  \defaultfontfeatures{Scale=MatchLowercase}
  \defaultfontfeatures[\rmfamily]{Ligatures=TeX,Scale=1}
\fi
\usepackage{lmodern}
\ifPDFTeX\else
  % xetex/luatex font selection
\fi
% Use upquote if available, for straight quotes in verbatim environments
\IfFileExists{upquote.sty}{\usepackage{upquote}}{}
\IfFileExists{microtype.sty}{% use microtype if available
  \usepackage[]{microtype}
  \UseMicrotypeSet[protrusion]{basicmath} % disable protrusion for tt fonts
}{}
\makeatletter
\@ifundefined{KOMAClassName}{% if non-KOMA class
  \IfFileExists{parskip.sty}{%
    \usepackage{parskip}
  }{% else
    \setlength{\parindent}{0pt}
    \setlength{\parskip}{6pt plus 2pt minus 1pt}}
}{% if KOMA class
  \KOMAoptions{parskip=half}}
\makeatother
\usepackage{xcolor}
\usepackage[margin=21mm]{geometry}
\usepackage{longtable,booktabs,array}
\usepackage{calc} % for calculating minipage widths
% Correct order of tables after \paragraph or \subparagraph
\usepackage{etoolbox}
\makeatletter
\patchcmd\longtable{\par}{\if@noskipsec\mbox{}\fi\par}{}{}
\makeatother
% Allow footnotes in longtable head/foot
\IfFileExists{footnotehyper.sty}{\usepackage{footnotehyper}}{\usepackage{footnote}}
\makesavenoteenv{longtable}
\setlength{\emergencystretch}{3em} % prevent overfull lines
\providecommand{\tightlist}{%
  \setlength{\itemsep}{0pt}\setlength{\parskip}{0pt}}
\setcounter{secnumdepth}{-\maxdimen} % remove section numbering
\ifLuaTeX
  \usepackage{selnolig}  % disable illegal ligatures
\fi
\IfFileExists{bookmark.sty}{\usepackage{bookmark}}{\usepackage{hyperref}}
\IfFileExists{xurl.sty}{\usepackage{xurl}}{} % add URL line breaks if available
\urlstyle{same}
\hypersetup{
  pdftitle={Seven unresolved fusion rings: audit and theoretical continuation},
  colorlinks=true,
  linkcolor={blue},
  filecolor={Maroon},
  citecolor={Blue},
  urlcolor={blue},
  pdfcreator={LaTeX via pandoc}}

\title{Seven unresolved fusion rings: audit and theoretical
continuation}
\usepackage{etoolbox}
\makeatletter
\providecommand{\subtitle}[1]{% add subtitle to \maketitle
  \apptocmd{\@title}{\par {\large #1 \par}}{}{}
}
\makeatother
\subtitle{The three previous exclusions are retained; seven problems
remain open}
\author{}
\date{27 October 2026}

\begin{document}
\maketitle

\hypertarget{outcome}{%
\subsection{Outcome}\label{outcome}}

The fresh review retains the general-complex exclusions of
\textbf{R9\_2\_1607, R9\_2\_1975, and R9\_2\_1976}. It found no gap
requiring their withdrawal. The inherited spherical/Galois and Morita
reductions remain explicit dependencies of these conclusions.

This continuation gives \textbf{no additional complete census-ring
exclusion and no new census-ring model}. All seven requested rings
remain unresolved. Each was investigated separately; the new results
below include categorical restrictions and exclusions of construction
routes, which must not be counted as exclusions of the rings themselves.

\begin{longtable}[]{@{}
  >{\raggedright\arraybackslash}p{(\columnwidth - 4\tabcolsep) * \real{0.1900}}
  >{\raggedright\arraybackslash}p{(\columnwidth - 4\tabcolsep) * \real{0.6500}}
  >{\raggedright\arraybackslash}p{(\columnwidth - 4\tabcolsep) * \real{0.1600}}@{}}
\toprule\noalign{}
\begin{minipage}[b]{\linewidth}\raggedright
Ring
\end{minipage} & \begin{minipage}[b]{\linewidth}\raggedright
Principal additional result
\end{minipage} & \begin{minipage}[b]{\linewidth}\raggedright
Census status
\end{minipage} \\
\midrule\noalign{}
\endhead
\bottomrule\noalign{}
\endlastfoot
R6\_4\_155 & Nonpointed algebra family reduced to two cases; one algebra
is forced but returns to a known module category. & Unresolved \\
R6\_8\_1660 & One center subcase removed; a pairing obstruction also
tests proposed Morita realizations. & Unresolved \\
R8\_3\_928 & Its nonpointed modular centralizer must be simple, with
only one nontrivial connected étale algebra. & Unresolved \\
R8\_3\_941 & Its rank-22 fermion centralizer must be nonsplit and have
no proper nonpointed fusion subcategory. & Unresolved \\
R8\_3\_1154 & All six third indicators are forced to be 2; all six
trivalent rotations are identity. & Unresolved \\
R8\_3\_1768 & A pairing obstruction rules out the higher-Haagerup Morita
construction family. & Unresolved \\
R9\_2\_137 & A nonpointed algebra is forced, but returns to an existing
Morita route; both TY signs pass the strengthened exact test. &
Unresolved \\
\end{longtable}

The source snapshot is the classification repository at
\texttt{6a426acf6465dd674351be92dcca41b72f371521}. The publication
applies the three separately audited rank-nine exclusions. The
seven-case continuation adds restrictions without changing those seven
statuses. The computations in this pass were short exact replays or
bounded certificate-discovery tests, not long associator searches.

\hypertarget{fresh-audit-of-the-three-exclusions}{%
\subsection{1. Fresh audit of the three
exclusions}\label{fresh-audit-of-the-three-exclusions}}

For \textbf{R9\_2\_1975}, the review checked the opposite-bimodule and
normalizer conventions, cancellation for arbitrary complex cocycles, the
strict even-\(V_4\) subgroup used in the projective parity argument, and
the normalized character formula for the distinguished Galois central
object. The required restriction remains

\[
F(B)=1+X_0+X_1+X_2+X_3.
\]

In the rotation branch its multiplicity at \(X_0\) is odd, whereas the
forced anticommuting stabilizer operators require it to be even.
Reflection shifts are separately impossible. The audit preserves the
eighth-root coboundary witness: a fourth-root-only test would be
invalid. The inherited extreme-branch argument and complete signed gate
were also checked as stated dependencies. The fresh gate replay retained
two branches and certified completeness.

For \textbf{R9\_2\_1607 and R9\_2\_1976}, the review specifically
checked that pivotal symmetry of the Frobenius pairing, rather than
merely braided commutativity, is available. The selfdual
multiplicity-one Galois summand of the canonical algebra has a symmetric
nondegenerate pairing. Transporting that pairing forces all four large
second indicators in the zero-shift Morita category to be \(+1\). The
dimension-weighted second-indicator total is Morita invariant.
Comparison of irrational coefficients then forces the same sum to be
four in the original category. The original rings have, respectively,
only two and zero selfdual large simples. This is the contradiction. The
irrational-coefficient comparison avoids a dependence on TY sign
conventions. Both complete signed gates were replayed successfully.

These arguments do not assume that arbitrary tensor functors preserve
numerical indicators. Their pairing transport and Morita-invariant total
are different, explicitly justified steps. Full fresh reviews are
\texttt{audit\_parity/AUDIT\_R1975.md} and \texttt{audit\_fs2/AUDIT.md}.

\hypertarget{a-useful-common-obstruction-pairings-under-morita-forgetful-functors}{%
\subsection{2. A useful common obstruction: pairings under Morita
forgetful
functors}\label{a-useful-common-obstruction-pairings-under-morita-forgetful-functors}}

Let \(\mathcal C\) be pseudo-unitary, and let \(A\) be a connected
separable algebra. With its positive module trace, \(A\) can be equipped
with the compatible symmetric separable Frobenius structure. The functor

\[
U:{}_A\mathcal C_A\longrightarrow\mathcal C
\]

is pivotal Frobenius monoidal. It preserves a chosen nondegenerate
pairing and its pivotal symmetry sign. In particular, if a selfdual
simple bimodule \(Q\) has second indicator \(\varepsilon\), then every
selfdual simple constituent of \(U(Q)\) having the opposite indicator
must occur with even multiplicity.

Indeed, on such an isotypic block the multiplicity space must carry a
nondegenerate alternating form. Its dimension is even. This supplies a
short necessary test for proposed Morita restrictions. It does not
assert equality of the numerical indicators of \(Q\) and \(U(Q)\).

The positive pivotal normalization matters. The bimodule dimension is
\(\dim_{\mathcal C}(U(Q))/\dim_{\mathcal C}(A)\), so the induced
structure is the canonical positive one. The relevant primary inputs are
Schaumann, \emph{Traces on module categories over fusion categories},
Proposition 5.8 and Theorem 6.6, and Mulevičius, \emph{Condensation
inversion and Witt equivalence via generalised orbifolds}, Example 4.18.

\hypertarget{r6_4_155-forced-algebras-and-a-morita-loop}{%
\subsection{3. R6\_4\_155: forced algebras and a Morita
loop}\label{r6_4_155-forced-algebras-and-a-morita-loop}}

Write the invertibles as \(1,h,g,hg\), and the other simples as \(X,Y\),
with \(X^*=Y\), \(hX=X\), \(gX=Y\), and

\[
XY=1+h+4(X+Y),\qquad X^2=Y^2=g+hg+4(X+Y).
\]

Put \(d=4+3\sqrt2\). Consider a connected separable algebra with
underlying object \(A_k=1+kX+kY\). Four free pointed modules are simple
and distinct. Subtracting them from the induction Gram matrix leaves
nonnegative integral vectors \(a,b\) with

\[
\|a\|^2=\|b\|^2=1+8k-4k^2,\qquad a\cdot b=8k-4k^2.
\]

Thus \(k\le2\). At \(k=1\) the only viable residual modules have
underlying objects \(2X+Y\) and \(X+2Y\). The other Gram possibilities
would contain a module with underlying object \(X\), whose internal end
would have dimension \(2+4\sqrt2/3\), not an algebraic integer. At
\(k=2\) the two residual modules have underlying objects \(X\) and
\(Y\). Each surviving case has a unique necessary based module of rank
six.

There is also a conditional existence statement. The canonical central
subalgebra corresponding to the pointed subcategory has underlying
central object \(1+B\), where \(B\) is the Galois-unit object. Its
forgetful algebra splits through its augmentation as

\[
F(1+B)\cong1\times A_2,\qquad [A_2]=1+2X+2Y.
\]

So \(A_2\) is forced in any hypothetical categorification. However the
residual module underlying \(X\) has internal end \(1+h\). This returns
to the already available pointed-algebra module category and supplies
neither a construction of the original category nor a contradiction. The
potentially different algebra \(A_1\) is not forced or constructed.

An initially tempting trace contradiction was rejected during review:
the trace on the four pointed free modules must be included. The total
is eight, not four. The accepted proof and replay are
\texttt{audit\_fs2/R155\_ALGEBRA\_REDUCTION.md} and its companion
verifier.

\hypertarget{r6_8_1660-a-narrower-valid-center-restriction}{%
\subsection{4. R6\_8\_1660: a narrower valid center
restriction}\label{r6_8_1660-a-narrower-valid-center-restriction}}

The two remaining possibilities for the negative charge-conjugation
representation meeting the four spin phases were

\[
(16,2)\oplus(20,2),\qquad (44,2).
\]

The second possibility cannot occur in low frame 0. In that frame all
high restriction columns have equal \(X\) and \(gX\) multiplicities. The
phase-12 Gram identities force exactly one column there. It is selfdual,
leaving no negative charge-conjugation line at that phase, whereas
\((44,2)\) requires one.

A proposed extension to all of frame 1 failed independent checking. The
total asymmetric-column budget is eight, not four; the earlier four
counted only one member of each pair. The remaining budget can support
phase-36 columns that change an apparent half-integral Verlinde entry to
the integer zero. Additional negative irreducibles away from the spin
phases also cannot be omitted. Thus the entire \((44,2)\) branch is not
excluded, and both overall alternatives remain.

The pairing test in Section 2 does exclude two specific proposed Morita
module realizations, labelled \texttt{III\_mixed0} and
\texttt{IV\_mixed1}. A selfdual simple of their respective duals would
restrict to \(\sigma+T\) and \(X+gX+\sigma+T\). Here \(\sigma\) has
indicator \(-1\) and every other displayed simple has indicator \(+1\).
Their odd multiplicities cannot carry a pairing with a single symmetry
sign. The conditional route to R941 passes this test and is retained.
More precisely, any connected separable algebra \(1\oplus X\) in a
hypothetical pseudo-unitary R1660 category forces the R941 dual: the
complete module list forces type III, the mixed signs are now excluded,
and the other equal-sign dual was already ruled out. The inverse algebra
in that context is also \(1\oplus Y\), with \(Y\) selfdual of the same
FP dimension. Neither algebra has been constructed. Fixing the unit
leaves one multiplication component to the unit and three to \(X\), but
their associativity still depends on the unknown ambient associator.

There is also a further rigidity condition: fusion by one distinguished
Galois-unit simple acts on the negative charge-conjugation space as a
signed permutation. Its eigenvalues there are \(\pm1\); in the
normalized dual-difference basis it is a symmetric integral matrix whose
square is identity. Hence each row has precisely one entry \(\pm1\).
This has not yet yielded a contradiction.

No forced algebra or categorical realization resolving R1660 was found.
The independent reviews are \texttt{audit\_fs2/R1660\_SPIN\_AUDIT.md}
and \texttt{audit\_fs2/R1660\_BIMODULE\_PARITY\_AUDIT.md}.

\hypertarget{r8_3_928-and-r8_3_941-simple-centralizers}{%
\subsection{5. R8\_3\_928 and R8\_3\_941: simple
centralizers}\label{r8_3_928-and-r8_3_941-simple-centralizers}}

Let \(D\) be their common global dimension. Their only proper nontrivial
fusion subcategory is pointed of rank two. Consequently any central
fusion subcategory containing a nonintegral object forgets dominantly
onto the whole source category and has dimension at least \(D\).
Equality would make the source braided, impossible for these
noncommutative rings.

For \textbf{R928}, write its center as \(S\boxtimes\mathcal B\), where
\(S\) is Semion and \(\dim\mathcal B=D^2/2\). The unit is the only
integral simple in \(\mathcal B\). A proper nontrivial subcategory and
its centralizer would both dominate the source. Their dimensions would
therefore have product at least \(D^2\), contradicting the centralizer
formula \(D^2/2\). Thus \(\mathcal B\) is \textbf{simple}, not just
indecomposable as a tensor product. Its possible ranks are 17, 20, and
23.

The earlier all-module rigidity theorem implies a further statement: the
entire center has exactly three connected étale algebra classes, forming
\(1\subset A\subset I(1)\), with \(\dim A=D/2\). The factor
\(\mathcal B\) contains exactly one nontrivial such algebra, namely
\(A\), and

\[
\mathcal B_A\simeq\mathcal C,\qquad
\mathcal B_A^0\simeq S^{\rm rev}.
\]

For \textbf{R941}, the fermion centralizer \(\mathcal E\) has rank 22,
dimension \(D^2/2\), and integral part exactly sVec. It cannot split as
\(\mathcal M\boxtimes\mathrm{sVec}\): that would make the center have
rank 33 or 44, whereas the complete retained list gives 34, 37, 40, or
43. The same dominant-image and centralizer argument then proves that
its only fusion subcategories are Vec, sVec, and \(\mathcal E\).

These results sharply specify the missing modular or slightly degenerate
categories, but do not prove them nonexistent. The current \emph{Anyon
Atlas} was checked; its higher-rank numerical search is not certified
exhaustive and cannot justify exclusion by absence. The full proofs and
exact rank checks are in \texttt{r8\_928\_941/REPORT.md}.

\hypertarget{r8_3_1154-every-trivalent-rotation-is-identity}{%
\subsection{6. R8\_3\_1154: every trivalent rotation is
identity}\label{r8_3_1154-every-trivalent-rotation-is-identity}}

The preceding actual equivariantization reduction gives a hypothetical
rank-ten category with a central sign object \(u\) and lifts \(x,ux\) of
\(X_4\), satisfying \(x^*=ux\). For odd \(n\), the two character
subspaces of \(\operatorname{Hom}(1,X_4^n)\) are the rotations on
\(\operatorname{Hom}(1,x^n)\) and \(\operatorname{Hom}(1,(ux)^n)\).
Spherical duality gives equal indicators of \(x\) and \(x^*\). Therefore

\[
\nu_n(X_4)=2\nu_n(x)\quad(n\text{ odd}).
\]

The same holds for \(X_5\). Their lift's cubic vacuum spaces are
one-dimensional, so their third indicators lie in
\(\{2,2\zeta_3,2\zeta_3^2\}\). This is stronger than allowing an
arbitrary sum of two cube roots.

Applying these discrete possibilities to the three already retained
center patterns forces both indicators to be two. Duality, the inner
autoequivalence, and a second small discrete check then force

\[
\boxed{\nu_3(X_i)=2\quad
i=1,2,4,5,6,7.}
\]

Each of these six cubic spaces is two-dimensional. Since its rotation
has order three, trace two means that the rotation is the identity. All
33 excluded discrete branches have exact integer separating
certificates; the six retained branch records have exact rational
feasibility witnesses. An independent implementation checked them. No
center pattern is eliminated outright.

The theorem and solver-free replay are in
\texttt{r8\_1154/ODD\_ROTATION\_THEOREM.md}; the independent audit is in
\texttt{audit\_parity/AUDIT\_R1154\_ODD\_ROTATION.md}.

A second short argument proves that the integral part of its center is
Vec. An integral central simple could forget only to copies of \(1\) and
\(g\), but \(g\) is not central in the fusion ring, forcing its
coefficient to vanish. The full adjoint subring then makes every central
object supported on the unit trivial. This applies directly to arbitrary
complex categorifications. It rules out nontrivial finite-group
extensions and equivariantizations throughout the Morita class. In the
auxiliary rank-ten category, the integral center is exactly double
Semion; its unique nontrivial Tannakian subcategory only gives back the
already known descent.

Three further direct central-condensation families are ruled out:
\(\mathrm{Fib}^{\boxtimes6}\),
\(GH(C_4,1)\boxtimes\mathrm{Fib}^{\boxtimes3}\), and
\(GH(C_4,1)^{\boxtimes2}\), with arbitrary pointed factors. In the last
case, the required index-three subgroup contains a central \(C_4\); its
condensation would force the impossible rank-two rule \(x^2=1+4x\).
These are specified-family obstructions, not a classification of every
possible construction.

\hypertarget{r8_3_1768-a-uniform-higher-haagerup-obstruction}{%
\subsection{7. R8\_3\_1768: a uniform higher-Haagerup
obstruction}\label{r8_3_1768-a-uniform-higher-haagerup-obstruction}}

Put \(u=6+\sqrt{37}\). The positive source indicators are
\((1,-1,1,1,1,1,1,1)\), and its weighted second-indicator total is
\(6u\). In a category with \(GH(C_6,2)\) fusion rules, the six large
simples are selfdual of dimension \(u\). The pointed contribution to
that total is \(1+\nu_2(g^3)\). Morita invariance and irrationality
force

\[
\nu_2(g^3)=-1,\qquad
\nu_2(g^j\sigma)=+1\quad(0\le j<6)
\]

whenever such a category is Morita equivalent to R1768.

First consider the proposed algebra \(A=1+g^j\sigma\). The negative
invertible of the R1768 dual has underlying dimension \(1+u\) in the
higher-Haagerup category. Irrationality forces its underlying object to
contain one pointed and one large simple, each once. Its negative
pairing cannot restrict to the latter, whose second indicator is
positive. Section 2 gives a contradiction. This rules out the route for
\textbf{every associator}, not only the original arithmetic precursor.

A stronger use of the same parity principle applies on the R1768 side.
For a negative selfdual simple in a Morita dual, all six positive
nonpointed source constituents must occur evenly. Thus the \(\sqrt{37}\)
coefficient of its underlying dimension must be even. The complete
previously retained list has 21 possible module factors. In 17 of them,
some simple module has internal end with an odd such coefficient.
Choosing that simple as generator prevents its dual from having a
negative-indicator invertible. Each of the other four factors has a
dimension-one internal end, so its module category is regular.

It follows that a Morita representative with a negative-indicator
invertible must come from the regular module category, and hence have
rank eight. In particular \textbf{no category with \(GH(C_6,2)\) fusion
rules is Morita equivalent to it}. This eliminates the entire
higher-Haagerup Morita family, including other algebras in that family.
It does not eliminate an unknown R1768 category itself.

The obstruction covers every associator and every algebra in that Morita
family; changing its arithmetic precursor cannot evade it.

\hypertarget{r9_2_137-a-forced-algebra-does-not-supply-a-new-category}{%
\subsection{8. R9\_2\_137: a forced algebra does not supply a new
category}\label{r9_2_137-a-forced-algebra-does-not-supply-a-new-category}}

In its remaining zero-shift rank-twelve Morita category, let
\(S=\sum_qX_q\), \(d=4+3\sqrt2\), and \(z\) be the central involution of
the pointed \(D_8\). The canonical étale subalgebra associated to the
pointed subcategory is forced to be \(1+B\), where \(B\) is the
Galois-unit simple. Its forgetful algebra has augmentation and splits as

\[
F(1+B)\cong1\times A,\qquad [A]=1+S.
\]

Thus this nonpointed algebra is forced, not merely a speculative
construction. Its induction Gram, however, gives eight modules
underlying \(g+S\) and four underlying \(X_q\). Each of the latter has
internal end \(1+z\). The module category is therefore the one already
obtained from the central-involution algebra. This is a Morita loop, not
a realization of a presently hypothetical category.

The previously established positive second indicators were then inserted
into both remaining TY-sign center systems. Each has an exact
nonnegative rational witness, with support 102 and 97, respectively. The
complete short tests took under four seconds each; the saved witnesses
replay without an optimizer. They show that this combination of
constraints still does not exclude the ring.

The proof, source bindings, matrices, and witnesses are in
\texttt{audit\_parity/R137\_CONTINUATION.md} and its companion files.

\hypertarget{scope-of-the-remaining-problems}{%
\subsection{9. Scope of the remaining
problems}\label{scope-of-the-remaining-problems}}

No based module, modular-data candidate, precursor category, or feasible
indicator witness has been promoted to a model of a census ring. The
missing ingredients are still actual categorical constructions or
complete obstructions for the surviving cases. The report separates
accepted theorems, inherited premises, and rejected shortcuts so that
the next investigation can start from the stronger frontier without
repeating invalid inferences.

\hypertarget{primary-references-used-for-the-new-arguments}{%
\subsubsection{Primary references used for the new
arguments}\label{primary-references-used-for-the-new-arguments}}

\begin{itemize}
\tightlist
\item
  Davydov, Müger, Nikshych, Ostrik, \emph{The Witt group of
  non-degenerate braided fusion categories}:
  \url{https://arxiv.org/abs/1009.2117}.
\item
  Drinfeld, Gelaki, Nikshych, Ostrik, \emph{On braided fusion categories
  I}: \url{https://arxiv.org/abs/0906.0620}.
\item
  Schaumann, \emph{Traces on module categories over fusion categories}:
  \url{https://arxiv.org/abs/1206.5716}.
\item
  Mulevičius, \emph{Condensation inversion and Witt equivalence via
  generalised orbifolds}: \url{https://arxiv.org/abs/2206.02611}.
\item
  Ng, Schauenburg, \emph{Frobenius-Schur indicators and exponents of
  spherical categories}: \url{https://arxiv.org/abs/math/0601012}.
\end{itemize}

\end{document}
